\section{Facilitation and Convergent Group Decision-Making (RQ1)}
\label{sec:facilitation-and-group-decisions}

Rechreche was sind eingelich die aufgaben eines Facilitators? was macht ein facilitator?

Welche rolle spielen collaborations tools wie white boards (Google Jamboard, Miro, etc.) in der facilitation?

% % ---------------------------------------------------------------------------------------------------------------------
% \subsection{The facilitator role}
% The book *Facilitation Fundamentals* by AJ\&Smart \cite{aj&smartFacilitationFundamentalsWorkbook}, one of the
% leading companies in the field of facilitation and facilitator training, introduces the 4C framework. This framework
% classifies methods that can be used in meetings or workshops into four key categories:

% \begin{itemize}
%   \item Collect: Suitable for systematically gathering and visualizing information within a team.
%   \item Choose: Narrows down a large number of ideas or content to a few selected elements.
%   \item Create: Uses design thinking methods to develop new approaches to solutions.
%   \item Commit: Ensures the team’s binding agreement on the decisions made and the results achieved.
% \end{itemize}

% As outlined in the chapter “Convergent and Divergent Thinking,” a clear distinction between the divergent and
% convergent phases is crucial. The 4C framework, which many facilitators rely on, enables
% the implementation of these best practices.

% \section{Convergent Decision-Making in Groups (RQ1)}

% ---------------------------------------------------------------------------------------------------------------------
\subsection{Group Support Systems and collaboration tools}
\label{sec:gss-collaboration-tools}

Recherechen zu Group Support systeme begannen in den 1980er jahren mit den stetigen herausfoderungen von
komplexerwerten problmen die in organsiation oft in gruppen / Teams gelöst werden müssen.
\cite[p. 628]{ackermannGroupSupportSystems2021}

Laut Akerman waren die Anfänge in diesem Forschungsbereich durch eine starke Fragmentierung geprägt.
Verschiedene Forschungsteams arbeiteten unabhängig voneinander und ohne Abstimmung an überlappenden Themen.
Dies führte dazu, dass es kein standardisiertes Vorgehen gab und sich unterschiedliche Bezeichnungen für
dasselbe Forschungsgebiet etablierten, wie etwa „Meeting Support Systems“, „Computer-Supported Cooperative
Work Systems“, „Group Support Systems“ oder „Group Decision Support Systems“ (GDSS).
\cite[p.628f]{ackermannGroupSupportSystems2021} Das übergeordnete Ziel besteht allerdgins darin, Meetings sowohl in
Bezug auf ihre Ergebnisqualität als auch auf ihre Prozessstruktur zu optimieren.
\cite[p.629]{ackermannGroupSupportSystems2021}

Hierbei kommen oft computerunterstützte tools zum Einsatz. So definierte Huber
\cite[p. 195]{huberIssuesDesignGroup1984a}:

\begin{quote}
  "GDSS consists of a set of software, hardware, and language components and procedures that support a group
  of people engaged in a decision-related meeting."
\end{quote}

Klassische und initiale Group Support Systems (GSS) variieren stark in ihrer Komplexität und ihrer Funktionalität.

TODO: Quellen ergänzen
Einige Systeme ermöglichen es Teilnehmenden, simultan und anonym Textbeiträge einzugeben, um elektronisches
Brainstorming, Kriterienanalysen oder Priorisierungsabstimmungen durchzuführen. Ein Beispiel hierfür ist
GroupSystems (später als ThinkTank bekannt, vermutlich basierend auf Plexsys)
\cite[p.634]{ackermannGroupSupportSystems2021}.
Andere, strukturorientierte Ansätze setzen auf qualitative Methoden wie das Causal Mapping, um komplexe
Problemstellungen als visuelle Ursache-Wirkungs-Netzwerke abzubilden. Ein Vertreter dieser Kategorie ist der
Group Explorer (ehemals COPE bzw. Decision Explorer) [4][5].
Darüber hinaus nutzen einige Systeme quantitative Multi-Kriterien-Entscheidungsanalysen (MCDA), bei denen
Handlungsalternativen anhand mathematisch gewichteter Kriterien bewertet werden. Hierzu zählen etwa
MeetingWorks oder Decision Conferencing [6].
Ergänzend dazu strukturieren weitere Ansätze unübersichtliche Debatten mithilfe des Issue-Based Information
System (IBIS). Dabei werden Beiträge systematisch in Kernfragen, Lösungsoptionen sowie Pro- und
Contra-Argumente unterteilt, wie es etwa beim Dialogue Mapping der Fall ist [9][10].

Um diese unterschiedlichen Systemfunktionen in verlässliche,
wiederholbare Interaktionsmuster zu überführen, wurden schließlich **ThinkLets** als standardisierte
Moderationsbausteine entwickelt[11].

Moderne Tools wie Miro, Conceptboard oder Figma Jam, die als digitale Whiteboard-Lösungen klassifiziert
werden können, greifen einige dieser ursprünglichen Methoden auf und verbinden sie mit den Möglichkeiten
modernern Technologien das disitrubiteve zusammenarbeiten ermöglihct. Integirerte KI-basierter Systeme,
erweitern diese Tools durch smarte funktionen wie z.B. die automatische bundlung von Ideen durch KI zu Themenclustern
oder die Anreicherung von Ideen durch KI-generierte Vorschläge.

Zudem gibt es diverse  sammlungen an Übungen, Methodne und Tools inkluse anweisungen zur Druchführung und
Hilfestellungen worauf zu achten ist, welche für verschiende Gruppensituationen
und Herausfoderungen in meetings geeigent sind wie z.B.:
\begin{itemize}
  \item Liberating Structures (LS): A collection of methods and tools for facilitation.
  \item 9 Spaces: A framework for facilitation.
  \item SessionLab Libary: SessionLab ist ein Tools zum erstellen von Workshop agendas. Es gibt eine von der
    Community gepflegt Libary in der viele Methoden geteilt werden.
  \item TODO: add more and add Foodnotes for each item
\end{itemize}

% ---------------------------------------------------------------------------------------------------------------------
\subsection{The facilitator role}
\label{sec:facilitation}

Bereits in den Anfängen der GSS-Forschung (Group Support Systems) wiesen Beranek et al. darauf hin, dass GSS
allein die Effektivität einer Sitzung vermutlich nicht gewährleisten können. Vielmehr hängt der Erfolg
maßgeblich davon ab wie diese Systeme genutzt werden würden. \cite[p. 147]{beranekFacilitationGroupSupport1993}
Hierbei würde die Qualität der Gruppensitzung maßgeblich von dem Facilitators abhängen.
\cite[p.146]{beranekFacilitationGroupSupport1993}

Im deutschsprachigen Kontext wird der Begriff „Facilitator“ häufig mit „Moderator“ übersetzt. Diese
Gleichsetzung ist jedoch unzureichend, da die Rollen unterschiedliche Schwerpunkte aufweisen: Während ein
Moderator primär die Steuerung einer Sitzung übernimmt, zielt die Tätigkeit eines Facilitators auf die
prozessuale und inhaltliche Unterstützung der Gruppe ab. Die Rolle des Facilitators umfasst zwar moderierende
Elemente, geht jedoch darüber hinaus

According to Haney, researchers in the 1990s agreed that facilitation is a combination of
“(1)managing relationships between people, tasks and technology, (2) structuring tasks, and (3) contributing to
the effective accomplishment of the meeting outcomes.” \cite[p. 73]{hayneFacilitatorsPerspectiveMeetings}

In his book "The Skilled Facilitator", regarded as a standard work in the field, organizational
psychologist Roger Schwarz defines the term “Facilitator” as follows:
(\cite{schwarzSkilledFacilitatorComprehensive2016} as cited in
\cite{parsonsdialogueWhatMakesGood} \& \cite[p.9]{aj&smartFacilitationFundamentalsWorkbook}):

\begin{quote}
  "Group facilitation is a process in which a person who is acceptable to all members of a group,
  substantively neutral, and has no decision making authority, intervenes to help a group improve the way it
  identifies and solves problems, and makes decisions, in order to increase the group’s effectiveness."
\end{quote}

% ---------------------------------------------------------------------------------------------------------------------
\subsection{Divergent and Convergent Thinking in Meetings}
\label{sec:divergent-convergent}
Divergent thinking is the process of generating new ideas and concepts [TODO: QUELLE]. Convergent thinking is
the process of evaluating and selecting the best ideas from a set of options. [TODO: QUELLE]

According to Anne Manning, generating new ideas and evaluating them are two fundamentally different
processes. However, she says that in meetings, people often try to carry out both steps at the same time.
\cite[min 02:30]{professional&executivedevelopment-harvarddceConvergentVsDivergent2024}
Immediately evaluating and weighing an idea hinders the creative process of generating ideas.

A predominance of divergent thinking without sufficient integration of convergent thought processes can lead
to unreflective and excessive modifications. \cite{cropleyPraiseConvergentThinking2006} Corply
\cite{cropleyPraiseConvergentThinking2006} further argues that both a lack and an excess of
convergent or divergent thinking have adverse effects on creativity.

In the context of facilitation, the terms "Convergence Zone / Phase" and "Divergence Zone / Phase" are
widely recognized. A recurring framework in this domain is the group process map, developed by facilitator
Sam Kaner, commonly referred to as "Diamond of Participatory Decision-Making" or "Kaner’s Diamond".
\cite{kanerFacilitatorsGuideParticipatory2014}
Kaner’s Diamond model describes typical phases that groups go through in the decision-making process. The left
side represents the Divergent Zone, where diverse perspectives are expressed and disagreements are explored.
The right side embodies the Convergent Zone, where opinions align and the process moves toward closure. At
the center of the Diamond lies the "Groan Zone", a transitional phase where the tension between divergence
and convergence often manifests.

TODO: add figure of Kaner's Diamond

% ---------------------------------------------------------------------------------------------------------------------
\subsection{Biases in information exchange (Hidden profiles)}
\label{sec:hidden-profiles}

Diese Arbeit geht von der Hypothese aus, dass ein proaktives Meeting-Supporsystem insbesondere dort
eine Lücke schließen kann, wo die Gruppe ihre eigene Ineffizienz nicht erkennt. Ein interaktives
System setzt voraus, dass Teilnehmende ein Problem bemerken und um Unterstützung bitten. Genau das
entfällt, wenn der Prozess subjektiv reibungslos wirkt, obwohl er systematisch verzerrt ist.

Typischerweise liegt diese Unbewusstheit bei Biases vor, etwa Confirmation Bias, Groupthink,
Authority Bias, Halo-Effekt, Anchoring oder Status-Quo-Bias, um nur ein paar zu nennen.
Nicht alle diese Verzerrungen greifen jedoch am selben Punkt an.
Authority Bias oder Halo-Effekt betreffen vor allem, wessen Beitrag
oder welcher Gesamteindruck übergewichtet wird. Confirmation Bias und Groupthink betreffen
dagegen, welche Informationen überhaupt geteilt und diskutiert werden.

Für die konvergente Phase wird daher angenommen, dass zunächst vor allem der Informationsaustausch
als grundlage entscheidungsrelevant ist.

Dieses Problem wird in der Literatur als Hidden-Profile-Paradigma beschreiben. In Meetings verfügen Teilnehmende
typischerweise über unterschiedliche Informationen. Einen Teil kennen alle (geteilte Informationen),
einen anderen Teil nur einzelne Personen (ungeteilte/private Informationen). Ein Hidden Profile liegt vor,
wenn die bessere Lösung nur sichtbar wird, sobald die ungeteilten Informationen zusammengeführt
werden. Gruppen diskutieren jedoch vorwiegend geteilte Informationen und ungeteilte informationen kommen zu
krz oder werden nicht erwähnt.
\cite{stasserExpertRolesInformation1995,arkesInformationExchangeGroup2009,bakerEnhancingGroupDecision2010}.

Confirmation Bias und Groupthink verstärken diesen blinden Fleck, weil einzigartige Informationen den bestehenden
Konsens stören \cite{bakerEnhancingGroupDecision2010,schulz-hardtGroupDecisionMaking2006}.
Ohne externe Intervention bleibt die Ineffizienz unbemerkt.

Laut Verhein-Jarren et al. ist Voraussetzung für eine gute Kommunikation das Erkennen des eigenen
Persönlichkeitsstils sowie das Wissen darum, dass es verschiedene Bezugswelten gibt, in denen sich die
Beteiligten während des Gesprächs bewegen. Menschen neigen dazu, andere unbewusst nach eigenen Mustern zu
interpretieren. Laut Verhein-Jarren et al. sei ein Durchbrechen dieses Automatismus der Wahrnehmung
notwendig, um einen geteilten Bezugsrahmen zu schaffen. \cite[p.
9]{verhein-jarrenGespraechsfuehrungTechnischenBerufen2018}

Innerhlab des bleichen bezugsrahmen, kann es dann allerdgins druch aus auch zu meindugnsverscheidnheiten
kommen, was nicht unbedingt schlecht ist. Schulz-Hardt et al. (2006) zeigen empirisch, dass
Meinungsverschiedenheiten (Dissent) die Diskussionsintensität massiv erhöhen und die Lösungsrate des Hidden
Profiles drastisch steigern. \cite{schulz-hardtGroupDecisionMaking2006}

%%%%%%%%%%%%%%%%

Geteilte vs. Ungeteilte Informationen: Das Hidden Profil

Ein typisches Experiment, um Entscheidungsprozesse in Gruppen zu untersuchen sieht folgendermaßen aus: Ein
Team hat die Aufgabe aus unterschiedlichen BewerberInnen den Besten oder die Beste auszuwählen. Dazu erhält
jedes Team-Mitglied Informationen über die einzelnen BewerberInnen. Einen Teil der Information erhalten alle
Teammitglieder. Diese Informationen sind also im Team geteilte Informationen. Einen anderen Teil der
Informationen erhält jeweils nur ein Team-Mitglied. Diese Informationen sind also ungeteilt. Die Aufgabe ist
nun so konstruiert, dass sich das Team nur dann für den richtigen Bewerber entscheiden kann, wenn alle
Informationen – geteilte und ungeteilte – berücksichtigt werden. Werden lediglich die geteilten Informationen
berücksichtigt, trifft das Team eine falsche Entscheidung.

Eine Vielzahl von Studien mit dieser oder ähnlichen Aufgaben zeigt, dass Teams meist nicht in der Lage sind,
ein solches Hidden-Profile zu lösen: Geteilte Informationen werden ausgetauscht und diskutiert, ungeteilte
Informationen fallen unter den Tisch.
Erklärungen für das Phänomen Hidden Profil

Es gibt unterschiedliche Erklärungen dafür, warum in Gruppendiskussionen eher über geteilte Informationen
gesprochen wird:

Wahrscheinlichkeit: Mehr Personen verfügen über die gleiche geteilte Information. Deshalb ist die
Wahrscheinlichkeit höher, das geteilte Informationen in der Gruppe diskutiert werden, im Vergleich zu
Informationen, über die nur eine einzelne Person verfügt (ungeteilte Informationen).
Selbstwerterhaltung: Menschen achten darauf, ein positives Selbstwertgefühl zu erhalten. Wird eine geteilte
Information der Gruppe mitgeteilt, werden andere Personen in der Gruppe dieser Information zustimmen, bei
ungeteilten Informationen steigt die Gefahr, kritisiert zu werden. Das gefährdet ein positives
Selbstwertgefühl, weshalb eher ohnehin schon geteilte Informationen der Gruppe mitgeteilt werden.
Group Think: Jedes Gruppenmitglied hat das Bedürfnis, sich der Meinung der Gruppe anzuschließen. Die Gruppe
als Ganzes strebt nach Konsens. Das kann dazu führen, dass kritische oder widersprüchliche Informationen
nicht ausreichend berücksichtigt werden.

Lösungsstrategien für das Hidden Profil

Die Forschung zu Hidden-Profiles basiert auf eher künstlichen Gruppensituationen. Dennoch lassen sich aus den
Ergebnissen Strategien ableiten, Gruppenentscheidungen zu verbessen.

Effekt kennen: Wer die Anfälligkeit von Gruppen kennt, ungeteilte Informationen zu vernachlässigen, kann
gezielt gegensteuern. Eine Möglichkeit ist z. B. nachzufragen, ob alle Informationen berücksichtigt wurden,
oder ob es Informationen gibt, die nicht ins Bild passen.
Expertise transparent machen: In einer heterogenen Gruppe hilft es, die unterschiedlichen Fach- und
Wissensgebiete der einzelnen Mitglieder transparent zu machen. So wird klar, dass unterschiedliches Wissen
und unterschiedliche Meinungen zu erwarten sind.
Technische Unterstützung: Technische Werkzeuge, die Informationen visualisieren oder strukturieren, helfen
insbesondere bei komplexen Entscheidungen, alle Informationen zu berücksichtigen.
Erst suchen, dann bewerten: Eine weitere Strategie ist es, das Sammeln und das Bewerten von Informationen als
zwei Prozesse voneinander zu trennen. So wird verhindert, dass Informationen, die wichtig sein könnten,
direkt abgewertet werden.
Hierarchien vermeiden: Oft teilen insbesondere Mitglieder mit niedrigem Status ihr ungeteiltes Wissen nicht,
wohingegen geteiltes Wissen, dass Mitglieder mit hohem Status eingebracht haben, detailliert diskutiert wird.
Eine Lösungsstrategie ist deshalb, dass sich die Gruppenleitung oder Mitglieder mit hohem Status zunächst mit
einer eigenen Bewertung zurückhalten.

%Literaturhinweis: Stasser, G., \& Titus, W. (1985). Pooling of unshared information in group decision making:
%Biased information sampling during discussion. Journal of Personality and Social Psychology, 48, 1467–1478.

%Literaturnachweis: Stasser, G., \& Titus, W. (1985). Pooling of unshared information in group decision making:
%Biased information sampling during discussion. Journal of Personality and Social Psychology, 48, 1467–1478.

%Zitieren als: Moskaliuk, J. (2013). Warum Gruppen falsch entscheiden. wissens.blitz (95).
%https://wissensdialoge.de/hidden_profile
